% ============================================================
% FICHE MÉTHODE — STS BIOAL 2e ANNÉE
% Couples (moyenne ; écart type) et (médiane ; interquartile)
% Version enrichie de graphiques illustratifs
% ============================================================

\documentclass[11pt,a4paper]{article}

\usepackage[utf8]{inputenc}
\usepackage[T1]{fontenc}
\usepackage[french]{babel}
\usepackage{amsmath,amssymb}
\usepackage{geometry}
\usepackage{enumitem}
\usepackage{booktabs}
\usepackage{array}
\usepackage{tabularx}
\usepackage{xcolor}
\usepackage{tcolorbox}
\usepackage{fancyhdr}
\usepackage{multicol}
\usepackage{tikz}
\usepackage{pgfplots}
\pgfplotsset{compat=1.18}
\usetikzlibrary{babel}
\usetikzlibrary{arrows.meta,patterns}
\usepackage{lastpage}

\tcbuselibrary{skins,breakable}

\geometry{margin=2cm, headheight=15pt}

% ------------------------------------------------------------
% Couleurs
% ------------------------------------------------------------
\definecolor{primary}{RGB}{31,78,121}
\definecolor{soft}{RGB}{230,239,247}
\definecolor{amberbg}{RGB}{255,244,224}
\definecolor{amberborder}{RGB}{240,192,112}
\definecolor{greenbg}{RGB}{233,246,234}
\definecolor{greenborder}{RGB}{182,224,184}
\definecolor{redbg}{RGB}{253,236,236}
\definecolor{redborder}{RGB}{242,179,171}
\definecolor{purplebg}{RGB}{240,235,250}
\definecolor{purpleborder}{RGB}{190,170,220}
\definecolor{grayline}{RGB}{120,130,140}
\definecolor{accent}{RGB}{200,80,60}

% ------------------------------------------------------------
% En-tête / pied de page
% ------------------------------------------------------------
\pagestyle{fancy}
\fancyhf{}
\fancyhead[L]{\textcolor{primary}{\textbf{BTS BioALC — 2\textsuperscript{ème} année}}}
\fancyhead[R]{\textcolor{primary}{Fiche méthode — Statistiques descriptives}}
\fancyfoot[L]{\textcolor{grayline}{\small Nom : \dotfill}}
\fancyfoot[C]{\textcolor{grayline}{\small \thepage/\pageref{LastPage}}}
\fancyfoot[R]{\textcolor{grayline}{\small Couples (moyenne ; $\sigma$) et (médiane ; IQR)}}
\renewcommand{\headrulewidth}{0.4pt}
\renewcommand{\footrulewidth}{0.4pt}

% ------------------------------------------------------------
% Boîtes
% ------------------------------------------------------------
\newtcolorbox{consigne}[1][]{
  colback=soft, colframe=primary, boxrule=0.8pt, arc=4pt,
  fonttitle=\bfseries, coltitle=white, colbacktitle=primary,
  left=6pt, right=6pt, top=4pt, bottom=4pt,
  breakable, #1
}

\newtcolorbox{contexte}{
  colback=amberbg, colframe=amberborder, boxrule=0.8pt, arc=4pt,
  title={\textbf{Contexte professionnel}}, fonttitle=\bfseries,
  coltitle=black, colbacktitle=amberborder,
  left=6pt, right=6pt, top=4pt, bottom=4pt,
  breakable
}

\newtcolorbox{rappel}{
  colback=greenbg, colframe=greenborder, boxrule=0.8pt, arc=4pt,
  title={\textbf{Rappel — Formules}}, fonttitle=\bfseries,
  coltitle=black, colbacktitle=greenborder,
  left=6pt, right=6pt, top=4pt, bottom=4pt,
  breakable
}

\newtcolorbox{exemple}{
  colback=white, colframe=primary, boxrule=0.8pt, arc=4pt,
  title={\textbf{Exemple d'application — Calcul détaillé}}, fonttitle=\bfseries,
  coltitle=white, colbacktitle=primary,
  left=6pt, right=6pt, top=4pt, bottom=4pt,
  breakable
}

\newtcolorbox{interpretation}{
  colback=purplebg, colframe=purpleborder, boxrule=0.8pt, arc=4pt,
  title={\textbf{Interprétation}}, fonttitle=\bfseries,
  coltitle=black, colbacktitle=purpleborder,
  left=6pt, right=6pt, top=4pt, bottom=4pt,
  breakable
}

\newtcolorbox{attention}{
  colback=redbg, colframe=redborder, boxrule=0.8pt, arc=4pt,
  title={\textbf{Point de vigilance}}, fonttitle=\bfseries,
  coltitle=black, colbacktitle=redborder,
  left=6pt, right=6pt, top=4pt, bottom=4pt,
  breakable
}

\newtcolorbox{choix}{
  colback=soft, colframe=primary, boxrule=1pt, arc=4pt,
  title={\textbf{Critères de choix}}, fonttitle=\bfseries,
  coltitle=white, colbacktitle=primary,
  left=6pt, right=6pt, top=4pt, bottom=4pt,
  breakable
}

\newtcolorbox{graphbox}[1][]{
  colback=white, colframe=grayline, boxrule=0.6pt, arc=4pt,
  title={\textbf{Figure}}, fonttitle=\bfseries,
  coltitle=white, colbacktitle=grayline,
  left=4pt, right=4pt, top=2pt, bottom=2pt,
  breakable, #1
}

% ------------------------------------------------------------
% Lignes d'écriture
% ------------------------------------------------------------
\newcommand{\lignes}[1]{%
  \vspace{0.15cm}%
  \foreach \x in {1,...,#1}{%
    \noindent\rule{\linewidth}{0.4pt}\par\vspace{0.5cm}%
  }%
}

\title{}
\author{}
\date{}

% ============================================================
\begin{document}
% ============================================================

% ------------------------------------------------------------
% Page de titre
% ------------------------------------------------------------
\begin{center}
  {\Large \textbf{Fiche méthode — Statistique descriptive}}\\[0.3cm]
  {\large \textbf{Les couples d'indicateurs}}\\[0.2cm]
  {\large \textbf{(moyenne ; écart type)} et (médiane ; interquartile)}\\[0.3cm]
  {\large BTS Bioanalyses en Laboratoire de Contrôle — 2\textsuperscript{ème} année}
\end{center}

\vspace{0.4cm}

\begin{consigne}
\textbf{Objectifs de la fiche :}
\begin{itemize}[leftmargin=*]
  \item Savoir \textbf{déterminer} le couple (moyenne ; écart type) et le couple (médiane ; interquartile) sur une série de mesures.
  \item Comprendre le \textbf{sens} et l'\textbf{intérêt} de l'écart type et de l'écart interquartile.
  \item Savoir \textbf{choisir} le couple d'indicateurs le plus adapté à la situation analysée.
  \item Appliquer ces notions à des \textbf{exemples concrets} de laboratoire de biologie médicale.
\end{itemize}
\end{consigne}

\vspace{0.3cm}

\begin{contexte}
En biologie médicale, toute série de mesures (taux de glucose, diamètre de colonies, concentration d'un analyte, temps de coagulation, etc.) doit être \textbf{synthétisée} par des indicateurs pertinents. Un \textbf{indicateur de tendance centrale} (moyenne ou médiane) ne suffit jamais seul : il faut lui associer un \textbf{indicateur de dispersion} (écart type ou écart interquartile). Ces deux indicateurs forment un \textbf{couple} qui décrit à la fois le \textbf{centre} et la \textbf{dispersion} de la série. Le choix du couple dépend de la \textbf{nature de la distribution} et de la \textbf{présence ou non de valeurs aberrantes}.
\end{contexte}

\vspace{0.4cm}

% ============================================================
\section*{1. Le couple (moyenne ; écart type)}
% ============================================================

\subsection*{1.1. Définitions et formules}

\begin{rappel}
Pour une série de $n$ valeurs $x_1, x_2, \dots, x_n$ :

\vspace{0.2cm}
\textbf{Moyenne :}
\[
\bar{x} = \frac{1}{n}\sum_{i=1}^{n} x_i
\]

\vspace{0.2cm}
\textbf{Variance (de la population) :}
\[
\sigma^2 = \frac{1}{n}\sum_{i=1}^{n} (x_i - \bar{x})^2
\]

\vspace{0.2cm}
\textbf{Écart type :}
\[
\sigma = \sqrt{\sigma^2} = \sqrt{\frac{1}{n}\sum_{i=1}^{n} (x_i - \bar{x})^2}
\]

\vspace{0.2cm}
\textbf{Variance et écart type corrigés (échantillon, souvent notés $s^2$ et $s$) :}
\[
s^2 = \frac{1}{n-1}\sum_{i=1}^{n} (x_i - \bar{x})^2
\qquad
s = \sqrt{s^2}
\]
\textit{En STS BioAL, on utilise le plus souvent $s$ (échantillon) car les mesures ne concernent qu'un échantillon de la production.}
\end{rappel}

\subsection*{1.2. Sens et intérêt de l'écart type}

\begin{interpretation}
L'\textbf{écart type} mesure la \textbf{dispersion moyenne} des valeurs autour de la moyenne. Il s'exprime dans la \textbf{même unité} que la variable (g/L, mm, s, etc.).

\vspace{0.2cm}
\textbf{Interprétation concrète :}
\begin{itemize}[leftmargin=*]
  \item Un \textbf{écart type faible} signifie que les valeurs sont \textbf{regroupées} autour de la moyenne : la série est \textbf{homogène}, la mesure est \textbf{reproductible}.
  \item Un \textbf{écart type élevé} signifie que les valeurs sont \textbf{dispersées} : la série est \textbf{hétérogène}, la mesure est \textbf{peu précise} ou la population est \textbf{variée}.
\end{itemize}

\vspace{0.2cm}
\textbf{Intérêt en BIOAL :}
\begin{itemize}[leftmargin=*]
  \item Évaluer la \textbf{fidélité} d'une méthode de mesure (norme ISO 5725).
  \item Contrôler la \textbf{reproductibilité} d'un analyseur.
  \item Détecter une \textbf{perte de maîtrise} d'un processus (carte de contrôle).
  \item Calculer un \textbf{intervalle de confiance} autour de la moyenne.
\end{itemize}
\end{interpretation}

% ------------------------------------------------------------
% FIGURE 1 : Courbe de Gauss avec écarts types
% ------------------------------------------------------------
\begin{graphbox}
\textbf{Figure 1 — Sens de l'écart type sur une distribution symétrique (loi normale).}
\vspace{0.2cm}

\begin{center}
\begin{tikzpicture}
\begin{axis}[
  width=13.5cm, height=6.2cm,
  axis lines=left,
  xlabel={\small Valeur de la variable},
  ylabel={\small Fréquence},
  xtick={-3,-2,-1,0,1,2,3},
  xticklabels={$\bar{x}-3s$,$\bar{x}-2s$,$\bar{x}-s$,$\bar{x}$,$\bar{x}+s$,$\bar{x}+2s$,$\bar{x}+3s$},
  ytick=\empty,
  domain=-3.8:3.8,
  samples=150,
  smooth,
  enlargelimits=false,
  clip=true,
  axis line style={grayline},
  tick label style={font=\small},
]
% Aire sous la courbe
\addplot[fill=primary!15, draw=primary, thick] {exp(-x^2/2)} \closedcycle;
% Traits verticaux ±1s
\draw[primary, dashed, thick] (axis cs:-1,0) -- (axis cs:-1,0.607);
\draw[primary, dashed, thick] (axis cs:1,0) -- (axis cs:1,0.607);
% Traits verticaux ±2s
\draw[accent, dashed, thick] (axis cs:-2,0) -- (axis cs:-2,0.135);
\draw[accent, dashed, thick] (axis cs:2,0) -- (axis cs:2,0.135);
% Traits verticaux ±3s
\draw[grayline, dotted, thick] (axis cs:-3,0) -- (axis cs:-3,0.011);
\draw[grayline, dotted, thick] (axis cs:3,0) -- (axis cs:3,0.011);
% Annotations
\node[primary] at (axis cs:0,0.75) {\small $\approx 68\%$};
\node[accent] at (axis cs:0,0.40) {\small $\approx 95\%$};
\node[grayline] at (axis cs:0,0.18) {\small $\approx 99{,}7\%$};
% Ligne moyenne
\draw[primary, very thick] (axis cs:0,0) -- (axis cs:0,1.0);
\node[primary, above] at (axis cs:0,1.0) {\small $\bar{x}$};
\end{axis}
\end{tikzpicture}
\end{center}

\vspace{0.1cm}
{\small \textbf{Lecture :} dans une distribution symétrique (dite « normale »), environ 68\,\% des valeurs se situent dans l'intervalle $[\bar{x}-s\,;\,\bar{x}+s]$, 95\,\% dans $[\bar{x}-2s\,;\,\bar{x}+2s]$ et 99,7\,\% dans $[\bar{x}-3s\,;\,\bar{x}+3s]$. Plus $s$ est petit, plus la courbe est \textbf{resserrée} autour de la moyenne.}
\end{graphbox}

\vspace{0.3cm}

\subsection*{1.3. Exemple dynamique — Contrôle qualité d'un analyseur de glucose}

\begin{exemple}
Un technicien réalise \textbf{8 mesures répétées} du même échantillon de contrôle de glucose (en g/L) :

\begin{center}
\begin{tabular}{|c|c|c|c|c|c|c|c|c|}
\hline
Mesure & 1 & 2 & 3 & 4 & 5 & 6 & 7 & 8 \\
\hline
$x_i$ (g/L) & 1{,}02 & 0{,}98 & 1{,}01 & 0{,}99 & 1{,}00 & 1{,}03 & 0{,}97 & 1{,}00 \\
\hline
\end{tabular}
\end{center}

\vspace{0.2cm}
\textbf{Étape 1 — Calcul de la moyenne :}
\[
\bar{x} = \frac{1{,}02 + 0{,}98 + 1{,}01 + 0{,}99 + 1{,}00 + 1{,}03 + 0{,}97 + 1{,}00}{8}
= \frac{8{,}00}{8} = 1{,}00 \text{ g/L}
\]

\textbf{Étape 2 — Calcul des écarts à la moyenne :}

\begin{center}
\begin{tabular}{|c|c|c|}
\hline
$x_i$ & $x_i - \bar{x}$ & $(x_i - \bar{x})^2$ \\
\hline
1{,}02 & $+0{,}02$ & $0{,}0004$ \\
\hline
0{,}98 & $-0{,}02$ & $0{,}0004$ \\
\hline
1{,}01 & $+0{,}01$ & $0{,}0001$ \\
\hline
0{,}99 & $-0{,}01$ & $0{,}0001$ \\
\hline
1{,}00 & $0{,}00$ & $0{,}0000$ \\
\hline
1{,}03 & $+0{,}03$ & $0{,}0009$ \\
\hline
0{,}97 & $-0{,}03$ & $0{,}0009$ \\
\hline
1{,}00 & $0{,}00$ & $0{,}0000$ \\
\hline
\textbf{Somme} & & $\mathbf{0{,}0028}$ \\
\hline
\end{tabular}
\end{center}

\textbf{Étape 3 — Calcul de la variance corrigée :}
\[
s^2 = \frac{0{,}0028}{8-1} = \frac{0{,}0028}{7} \approx 0{,}0004
\]

\textbf{Étape 4 — Calcul de l'écart type :}
\[
s = \sqrt{0{,}0004} \approx 0{,}02 \text{ g/L}
\]

\textbf{Résultat :} le couple (moyenne ; écart type) est $\boxed{(\bar{x} \,;\, s) = (1{,}00 \,;\, 0{,}02)}$ g/L.
\end{exemple}

\begin{interpretation}
La moyenne des 8 mesures est de \textbf{1,00 g/L} et l'écart type de \textbf{0,02 g/L}. Les mesures s'écartent en moyenne de \textbf{0,02 g/L} autour de la valeur cible. Cet écart type \textbf{faible} indique une \textbf{bonne fidélité} de l'analyseur : la méthode est reproductible. Si l'écart type dépassait le seuil fixé par le laboratoire (par exemple 0,05 g/L), l'analyseur devrait être \textbf{recalibré} ou \textbf{maintenu}.
\end{interpretation}

% ============================================================
\newpage
% ============================================================

\section*{2. Le couple (médiane ; écart interquartile)}

\subsection*{2.1. Définitions et formules}

\begin{rappel}
Pour une série de $n$ valeurs \textbf{triées par ordre croissant} :

\vspace{0.2cm}
\textbf{Médiane ($Me$ ou $Q_2$) :} valeur qui partage la série en deux groupes de même effectif.
\begin{itemize}[leftmargin=*]
  \item $n$ impair : $Me = x_{(n+1)/2}$ (valeur centrale).
  \item $n$ pair : $Me = \dfrac{x_{n/2} + x_{n/2+1}}{2}$ (moyenne des deux valeurs centrales).
\end{itemize}

\vspace{0.2cm}
\textbf{Quartiles :}
\begin{itemize}[leftmargin=*]
  \item $Q_1$ : premier quartile (25\,\% des valeurs lui sont inférieures ou égales).
  \item $Q_2$ : deuxième quartile = médiane.
  \item $Q_3$ : troisième quartile (75\,\% des valeurs lui sont inférieures ou égales).
\end{itemize}

\vspace{0.2cm}
\textbf{Écart interquartile (IQR) :}
\[
\text{IQR} = Q_3 - Q_1
\]
\end{rappel}

\subsection*{2.2. Sens et intérêt de l'écart interquartile}

\begin{interpretation}
L'\textbf{écart interquartile} mesure la dispersion des \textbf{50\,\% centraux} de la série. Il s'exprime dans la \textbf{même unité} que la variable. Il \textbf{ignore volontairement} les 25\,\% de valeurs les plus faibles et les 25\,\% de valeurs les plus élevées : il est donc \textbf{robuste aux valeurs extrêmes}.

\vspace{0.2cm}
\textbf{Interprétation concrète :}
\begin{itemize}[leftmargin=*]
  \item Un \textbf{IQR faible} : les 50\,\% centraux des valeurs sont \textbf{resserrés} autour de la médiane → la population est \textbf{homogène} sur son cœur.
  \item Un \textbf{IQR élevé} : les 50\,\% centraux sont \textbf{étalés} → la population est \textbf{hétérogène} sur son cœur.
\end{itemize}

\vspace{0.2cm}
\textbf{Intérêt en BIOAL :}
\begin{itemize}[leftmargin=*]
  \item Décrire une série \textbf{asymétrique} ou contenant des \textbf{valeurs aberrantes} (ex. : un patient hors norme).
  \item Comparer des populations \textbf{sans être biaisé} par un petit nombre de valeurs extrêmes.
  \item Construire une \textbf{boîte à moustaches} (box plot) : $Q_1$, $Me$, $Q_3$ + moustaches à $1{,}5 \times \text{IQR}$.
  \item Évaluer la \textbf{dispersion réelle} d'une population hétérogène (ex. : cohorte de patients d'âges différents).
\end{itemize}
\end{interpretation}

% ------------------------------------------------------------
% FIGURE 2 : Boîte à moustaches avec quartiles
% ------------------------------------------------------------
\begin{graphbox}
\textbf{Figure 2 — Boîte à moustaches : lecture de la médiane et de l'écart interquartile.}
\vspace{0.2cm}

\begin{center}
\begin{tikzpicture}[scale=1.0]
  % Axe horizontal gradué
  \draw[thick, ->] (0,0) -- (13,0);
  \foreach \x/\lab in {1/2, 3/4, 5/6, 7/8, 9/10, 11/12} {
    \draw (\x,0) -- (\x,-0.15);
    \node[below, font=\small] at (\x,-0.15) {\lab};
  }
  \node[below, font=\small] at (6,-0.7) {Valeur de la variable (unités)};

  % Boîte à moustaches (Q1=3, Me=6, Q3=9)
  \draw[primary, very thick, fill=primary!15] (3,1.5) rectangle (9,2.5);
  \draw[primary, very thick] (3,1.2) -- (3,2.8);
  \draw[primary, very thick] (9,1.2) -- (9,2.8);
  % Médiane
  \draw[accent, very thick] (6,1.5) -- (6,2.5);
  % Moustaches (min=1, max=11)
  \draw[primary, thick] (1,2) -- (3,2);
  \draw[primary, thick] (9,2) -- (11,2);
  \draw[primary, thick] (1,1.7) -- (1,2.3);
  \draw[primary, thick] (11,1.7) -- (11,2.3);

  % Annotations
  \node[above, primary, font=\small] at (3,2.8) {$Q_1$};
  \node[above, accent, font=\small] at (6,2.5) {$Me = Q_2$};
  \node[above, primary, font=\small] at (9,2.8) {$Q_3$};
  \node[below, primary, font=\small] at (1,1.7) {min};
  \node[below, primary, font=\small] at (11,1.7) {max};

  % Flèche IQR
  \draw[<->, thick, accent] (3,3.6) -- (9,3.6);
  \node[above, accent, font=\small] at (6,3.6) {IQR $= Q_3 - Q_1$};

  % Flèche 50% centraux
  \node[font=\small, primary] at (6,4.3) {50\,\% centraux des valeurs};
  \draw[->, primary, thick] (6,4.1) -- (6,3.9);
\end{tikzpicture}
\end{center}

\vspace{0.1cm}
{\small \textbf{Lecture :} la boîte (rectangle) contient les 50\,\% centraux des valeurs. Sa largeur est l'écart interquartile. La médiane est le trait à l'intérieur de la boîte. Les moustaches s'étendent jusqu'aux valeurs extrêmes non aberrantes (à $1{,}5 \times \text{IQR}$ au maximum). Les points situés au-delà sont des \textbf{valeurs aberrantes}.}
\end{graphbox}

\vspace{0.3cm}

\subsection*{2.3. Exemple dynamique — Diamètre de colonies sur 11 boîtes de Petri}

\begin{exemple}
On mesure le diamètre de colonies (en mm) sur \textbf{11 boîtes de Petri} :

\begin{center}
\begin{tabular}{|c|c|c|c|c|c|c|c|c|c|c|c|}
\hline
Boîte & 1 & 2 & 3 & 4 & 5 & 6 & 7 & 8 & 9 & 10 & 11 \\
\hline
$x_i$ (mm) & 3{,}8 & 4{,}1 & 4{,}5 & 4{,}9 & 5{,}2 & 5{,}3 & 5{,}6 & 6{,}0 & 6{,}4 & 7{,}1 & 9{,}8 \\
\hline
\end{tabular}
\end{center}

\vspace{0.2cm}
\textbf{Étape 1 — Tri croissant :} les valeurs sont déjà triées (de 3,8 à 9,8 mm).

\textbf{Étape 2 — Détermination de la médiane :} $n = 11$ (impair), la médiane est la 6\textsuperscript{e} valeur :
\[
Me = x_{(11+1)/2} = x_6 = 5{,}3 \text{ mm}
\]

\textbf{Étape 3 — Détermination de $Q_1$ :} $Q_1$ est la valeur telle que 25\,\% des valeurs lui soient inférieures. Pour $n = 11$, on prend la valeur au rang $\lceil 0{,}25 \times 11 \rceil = 3$ :
\[
Q_1 = x_3 = 4{,}5 \text{ mm}
\]

\textbf{Étape 4 — Détermination de $Q_3$ :} $Q_3$ est la valeur telle que 75\,\% des valeurs lui soient inférieures. Rang : $\lceil 0{,}75 \times 11 \rceil = 9$ :
\[
Q_3 = x_9 = 6{,}4 \text{ mm}
\]

\textbf{Étape 5 — Calcul de l'écart interquartile :}
\[
\text{IQR} = Q_3 - Q_1 = 6{,}4 - 4{,}5 = 1{,}9 \text{ mm}
\]

\textbf{Résultat :} le couple (médiane ; interquartile) est $\boxed{(Me \,;\, \text{IQR}) = (5{,}3 \,;\, 1{,}9)}$ mm.
\end{exemple}

\begin{interpretation}
La médiane des diamètres est de \textbf{5,3 mm} : la moitié des boîtes ont un diamètre inférieur à 5,3 mm, l'autre moitié un diamètre supérieur. L'écart interquartile est de \textbf{1,9 mm} : les 50\,\% centraux des diamètres s'étalent sur 1,9 mm (de 4,5 à 6,4 mm).

\vspace{0.2cm}
La valeur \textbf{9,8 mm} (boîte n\textsuperscript{o}11) est \textbf{très éloignée} des autres. Si l'on avait utilisé la moyenne, celle-ci aurait été « tirée vers le haut » par cette valeur aberrante. La médiane et l'IQR, eux, \textbf{ne sont pas affectés} par cette valeur extrême : ils décrivent correctement le \textbf{cœur} de la série.
\end{interpretation}

% ------------------------------------------------------------
% FIGURE 3 : Boîte à moustaches de la série (avec valeur aberrante)
% ------------------------------------------------------------
\begin{graphbox}
\textbf{Figure 3 — Boîte à moustaches de la série de 11 boîtes de Petri.}
\vspace{0.2cm}

\begin{center}
\begin{tikzpicture}[scale=0.85]
  % Axe horizontal
  \draw[thick, ->] (0,0) -- (13,0);
  \foreach \x/\lab in {0.5/3, 2/4, 3.5/5, 5/6, 6.5/7, 8/8, 9.5/9, 11/10} {
    \draw (\x,0) -- (\x,-0.15);
    \node[below, font=\small] at (\x,-0.15) {\lab};
  }
  \node[below, font=\small] at (6.5,-0.7) {Diamètre (mm)};

  % Boîte : Q1=4.5, Me=5.3, Q3=6.4  (échelle 1 mm = 1.2 cm)
  % Position : x = 1.5 + (valeur - 3.8)*1.1
  % Q1=4.5 -> x=2.27 ; Me=5.3 -> x=3.15 ; Q3=6.4 -> x=4.36 ; min=3.8 -> x=1.5 ; max non aberrant=7.1 -> x=5.13 ; aberrant=9.8 -> x=8.10
  \draw[primary, very thick, fill=primary!15] (2.27,1.5) rectangle (4.36,2.5);
  \draw[primary, very thick] (2.27,1.2) -- (2.27,2.8);
  \draw[primary, very thick] (4.36,1.2) -- (4.36,2.8);
  \draw[accent, very thick] (3.15,1.5) -- (3.15,2.5);
  % Moustaches
  \draw[primary, thick] (1.5,2) -- (2.27,2);
  \draw[primary, thick] (4.36,2) -- (5.13,2);
  \draw[primary, thick] (1.5,1.7) -- (1.5,2.3);
  \draw[primary, thick] (5.13,1.7) -- (5.13,2.3);
  % Valeur aberrante
  \draw[accent, very thick] (8.10,1.6) -- (8.10,2.4);
  \node[accent, above, font=\small] at (8.10,2.4) {9{,}8 (aberrante)};

  % Annotations
  \node[above, primary, font=\small] at (2.27,2.8) {$Q_1=4{,}5$};
  \node[above, accent, font=\small] at (3.15,2.5) {$Me=5{,}3$};
  \node[above, primary, font=\small] at (4.36,2.8) {$Q_3=6{,}4$};
  \node[below, primary, font=\small] at (1.5,1.7) {3{,}8};
  \node[below, primary, font=\small] at (5.13,1.7) {7{,}1};

  % Flèche IQR
  \draw[<->, thick, accent] (2.27,3.5) -- (4.36,3.5);
  \node[above, accent, font=\small] at (3.31,3.5) {IQR $=1{,}9$ mm};
\end{tikzpicture}
\end{center}

\vspace{0.1cm}
{\small \textbf{Lecture :} la valeur 9,8 mm est nettement détachée de la moustache droite (qui s'arrête à 7,1 mm) : elle est considérée comme \textbf{aberrante}. La médiane et l'IQR décrivent le cœur de la série sans être perturbés par cette valeur.}
\end{graphbox}

% ============================================================
\newpage
% ============================================================

\section*{3. Comparaison des deux couples d'indicateurs}

\begin{center}
\renewcommand{\arraystretch}{1.5}
\begin{tabularx}{\linewidth}{|p{4cm}|X|X|}
\hline
\textbf{Critère} & \textbf{Couple (moyenne ; écart type)} & \textbf{Couple (médiane ; IQR)} \\
\hline
\textbf{Centre} & Moyenne $\bar{x}$ & Médiane $Me$ \\
\hline
\textbf{Dispersion} & Écart type $s$ & Écart interquartile $Q_3 - Q_1$ \\
\hline
\textbf{Sensibilité aux valeurs extrêmes} & \textbf{Forte} : une seule valeur aberrante modifie $\bar{x}$ et $s$. & \textbf{Faible} (robuste) : les 25\,\% extrêmes sont ignorés. \\
\hline
\textbf{Type de distribution adapté} & Distribution \textbf{symétrique}, sans valeurs aberrantes, proche d'une loi normale. & Distribution \textbf{asymétrique}, avec valeurs aberrantes, ou effectif faible. \\
\hline
\textbf{Utilisation en BIOAL} & Fidélité d'une méthode, cartes de contrôle, intervalles de confiance, comparaison de moyennes (tests statistiques). & Description d'une population hétérogène, boîtes à moustaches, comparaison de cohortes, détection visuelle d'anomalies. \\
\hline
\textbf{Unité} & Même unité que la variable (g/L, mm, s\ldots). & Même unité que la variable. \\
\hline
\textbf{Information fournie} & Dispersion \textbf{moyenne} autour du centre. & Dispersion des \textbf{50\,\% centraux}. \\
\hline
\end{tabularx}
\end{center}

% ------------------------------------------------------------
% FIGURE 4 : Effet d'une valeur aberrante sur moyenne vs médiane
% ------------------------------------------------------------
\begin{graphbox}
\textbf{Figure 4 — Effet d'une valeur aberrante sur la moyenne et sur la médiane.}
\vspace{0.2cm}

\begin{center}
\begin{tikzpicture}[scale=0.9]
  % --- Série A : sans valeur aberrante ---
  \node[font=\small\bfseries, primary] at (0,4.2) {Série A — sans valeur aberrante};
  \draw[thick, ->] (0,0) -- (12,0);
  \foreach \x in {1,2,...,11} {
    \draw (\x,0) -- (\x,-0.1);
  }
  % Points
  \foreach \x in {3,4,4.5,5,5.5,6,6.5,7,7.5,8,9} {
    \fill[primary] (\x,1) circle (3pt);
  }
  % Moyenne et médiane
  \draw[accent, very thick] (6,0) -- (6,3);
  \node[accent, above, font=\small] at (6,3) {$\bar{x} \approx Me$};

  % --- Série B : avec valeur aberrante ---
  \node[font=\small\bfseries, primary] at (0,-2.5) {Série B — avec une valeur aberrante (à droite)};
  \draw[thick, ->] (0,-6) -- (12,-6);
  \foreach \x in {1,2,...,11} {
    \draw (\x,-6) -- (\x,-6.1);
  }
  % Points (mêmes valeurs + une valeur aberrante)
  \foreach \x in {3,4,4.5,5,5.5,6,6.5,7,7.5,8} {
    \fill[primary] (\x,-5) circle (3pt);
  }
  \fill[accent] (11.5,-5) circle (4pt);
  \node[accent, above, font=\small] at (11.5,-4.7) {aberrante};
  % Moyenne décalée
  \draw[accent, very thick] (7.2,-6) -- (7.2,-3);
  \node[accent, above, font=\small] at (7.2,-3) {$\bar{x}$ décalée};
  % Médiane quasi inchangée
  \draw[primary, very thick] (6.25,-6) -- (6.25,-3);
  \node[primary, above, font=\small] at (6.25,-2.5) {$Me$ stable};

  % Légende
  \node[primary, font=\small] at (13.5,-2) {\textbf{---} Médiane};
  \node[accent, font=\small] at (13.5,-3) {\textbf{---} Moyenne};
\end{tikzpicture}
\end{center}

\vspace{0.1cm}
{\small \textbf{Lecture :} dans la série A (distribution symétrique), moyenne et médiane sont confondues. Dans la série B, l'ajout d'une valeur aberrante à droite \textbf{déplace fortement la moyenne} vers la droite, alors que la \textbf{médiane reste quasiment stable}. C'est la raison pour laquelle on préfère le couple (médiane ; IQR) en présence de valeurs aberrantes.}
\end{graphbox}

\vspace{0.3cm}

\section*{4. Critères de choix entre les deux couples}

\begin{choix}
\textbf{Choisir le couple (moyenne ; écart type) lorsque :}
\begin{itemize}[leftmargin=*]
  \item La distribution est \textbf{symétrique} (ou proche de la loi normale).
  \item Il n'y a \textbf{pas de valeurs aberrantes} identifiées.
  \item On souhaite \textbf{comparer des moyennes} entre plusieurs séries (tests statistiques : test de Student, ANOVA).
  \item On veut évaluer la \textbf{fidélité} d'une méthode (norme ISO 5725).
  \item On construit une \textbf{carte de contrôle} (limites $\bar{x} \pm 2s$ ou $\bar{x} \pm 3s$).
\end{itemize}

\vspace{0.3cm}
\textbf{Choisir le couple (médiane ; IQR) lorsque :}
\begin{itemize}[leftmargin=*]
  \item La distribution est \textbf{asymétrique} (ex. : délais d'attente, concentrations biologiques).
  \item La série contient des \textbf{valeurs aberrantes} (patients hors norme, incident technique).
  \item L'effectif est \textbf{faible} ($n < 30$).
  \item On souhaite décrire une \textbf{population hétérogène} sans être biaisé par ses extrêmes.
  \item On construit une \textbf{boîte à moustaches}.
\end{itemize}

\vspace{0.3cm}
\textbf{Règle pratique :} en l'absence d'information sur la distribution, il est conseillé de \textbf{calculer les deux couples} et de les comparer. Si $\bar{x}$ et $Me$ sont \textbf{proches}, la distribution est probablement symétrique → on utilise (moyenne ; écart type). Si $\bar{x}$ et $Me$ sont \textbf{très différents}, la distribution est asymétrique ou contient des valeurs aberrantes → on utilise (médiane ; IQR).
\end{choix}

% ============================================================
\newpage
% ============================================================

\section*{5. Application comparée — Deux séries de mesures en laboratoire}

\begin{exemple}
\textbf{Situation.} Le laboratoire compare deux séries de mesures de glucose (en g/L) réalisées par deux analyseurs différents, sur 10 échantillons de contrôle.

\vspace{0.2cm}
\begin{center}
\begin{tabular}{|c|c|c|c|c|c|c|c|c|c|c|}
\hline
\textbf{Analyseur A} & 0{,}95 & 0{,}98 & 1{,}00 & 1{,}02 & 1{,}01 & 0{,}99 & 1{,}03 & 0{,}97 & 1{,}00 & 1{,}05 \\
\hline
\textbf{Analyseur B} & 0{,}80 & 0{,}90 & 0{,}95 & 0{,}98 & 1{,}00 & 1{,}02 & 1{,}05 & 1{,}10 & 1{,}20 & 1{,}30 \\
\hline
\end{tabular}
\end{center}

\vspace{0.2cm}
\textbf{Calculs — Analyseur A :}
\begin{align*}
\bar{x}_A &= \frac{0{,}95 + 0{,}98 + 1{,}00 + 1{,}02 + 1{,}01 + 0{,}99 + 1{,}03 + 0{,}97 + 1{,}00 + 1{,}05}{10} \\
&= \frac{10{,}00}{10} = 1{,}00 \text{ g/L}
\end{align*}
\[
s_A \approx 0{,}03 \text{ g/L}
\]
Médiane : $n = 10$ (pair), les 5\textsuperscript{e} et 6\textsuperscript{e} valeurs triées sont 0,99 et 1,00, donc $Me_A = 0{,}995$ g/L.
$Q_1 = 0{,}975$ g/L ; $Q_3 = 1{,}015$ g/L ; $\text{IQR}_A = 0{,}040$ g/L.

\vspace{0.2cm}
\textbf{Calculs — Analyseur B :}
\begin{align*}
\bar{x}_B &= \frac{0{,}80 + 0{,}90 + 0{,}95 + 0{,}98 + 1{,}00 + 1{,}02 + 1{,}05 + 1{,}10 + 1{,}20 + 1{,}30}{10} \\
&= \frac{10{,}30}{10} = 1{,}03 \text{ g/L}
\end{align*}
\[
s_B \approx 0{,}15 \text{ g/L}
\]
Médiane : 5\textsuperscript{e} et 6\textsuperscript{e} valeurs triées = 1,00 et 1,02, donc $Me_B = 1{,}01$ g/L.
$Q_1 = 0{,}95$ g/L ; $Q_3 = 1{,}10$ g/L ; $\text{IQR}_B = 0{,}15$ g/L.

\vspace{0.3cm}
\textbf{Synthèse :}

\begin{center}
\begin{tabular}{|l|c|c|}
\hline
& \textbf{Analyseur A} & \textbf{Analyseur B} \\
\hline
Moyenne $\bar{x}$ & 1,00 g/L & 1,03 g/L \\
\hline
Écart type $s$ & 0,03 g/L & 0,15 g/L \\
\hline
Médiane $Me$ & 0,995 g/L & 1,01 g/L \\
\hline
IQR & 0,040 g/L & 0,15 g/L \\
\hline
\end{tabular}
\end{center}
\end{exemple}

% ------------------------------------------------------------
% FIGURE 5 : Boîtes à moustaches comparées A et B
% ------------------------------------------------------------
\begin{graphbox}
\textbf{Figure 5 — Boîtes à moustaches comparées des analyseurs A et B.}
\vspace{0.2cm}

\begin{center}
\begin{tikzpicture}[scale=1.1]
  % Axe horizontal (glucose en g/L, de 0.75 à 1.35)
  \draw[thick, ->] (0,0) -- (13,0);
  \foreach \x/\lab in {0.5/0.80, 2/0.90, 3.5/1.00, 5/1.10, 6.5/1.20, 8/1.30} {
    \draw (\x,0) -- (\x,-0.15);
    \node[below, font=\small] at (\x,-0.15) {\lab};
  }
  \node[below, font=\small] at (6.5,-0.7) {Glucose (g/L)};

  % Conversion : 1 g/L = 20 cm ; x = (valeur - 0.75)*20
  % Analyseur A : Q1=0.975 -> x=4.5 ; Me=0.995 -> x=4.9 ; Q3=1.015 -> x=5.3
  % min=0.95 -> x=4.0 ; max=1.05 -> x=6.0
  \draw[primary, very thick, fill=primary!20] (4.5,2.5) rectangle (5.3,3.5);
  \draw[primary, very thick] (4.5,2.2) -- (4.5,3.8);
  \draw[primary, very thick] (5.3,2.2) -- (5.3,3.8);
  \draw[accent, very thick] (4.9,2.5) -- (4.9,3.5);
  \draw[primary, thick] (4.0,3) -- (4.5,3);
  \draw[primary, thick] (5.3,3) -- (6.0,3);
  \draw[primary, thick] (4.0,2.7) -- (4.0,3.3);
  \draw[primary, thick] (6.0,2.7) -- (6.0,3.3);
  \node[primary, font=\small\bfseries] at (1.5,3) {Analyseur A};
  \node[primary, font=\small] at (5.0,4.2) {IQR $=0{,}04$ g/L};

  % Analyseur B : Q1=0.95 -> x=4.0 ; Me=1.01 -> x=5.2 ; Q3=1.10 -> x=7.0
  % min=0.80 -> x=1.0 ; max=1.30 -> x=11.0
  \draw[primary, very thick, fill=primary!20] (4.0,0.5) rectangle (7.0,1.5);
  \draw[primary, very thick] (4.0,0.2) -- (4.0,1.8);
  \draw[primary, very thick] (7.0,0.2) -- (7.0,1.8);
  \draw[accent, very thick] (5.2,0.5) -- (5.2,1.5);
  \draw[primary, thick] (1.0,1) -- (4.0,1);
  \draw[primary, thick] (7.0,1) -- (11.0,1);
  \draw[primary, thick] (1.0,0.7) -- (1.0,1.3);
  \draw[primary, thick] (11.0,0.7) -- (11.0,1.3);
  \node[primary, font=\small\bfseries] at (1.5,1) {Analyseur B};
  \node[primary, font=\small] at (6.5,2.2) {IQR $=0{,}15$ g/L};
\end{tikzpicture}
\end{center}

\vspace{0.1cm}
{\small \textbf{Lecture :} la boîte de l'analyseur A est \textbf{étroite} (IQR faible), signe d'une méthode \textbf{fidèle}. La boîte de l'analyseur B est \textbf{large} (IQR élevé) et ses moustaches sont \textbf{très étendues}, signe d'une méthode \textbf{peu reproductible}. Les médianes (traits rouges) sont proches, mais la dispersion diffère fortement.}
\end{graphbox}

\vspace{0.3cm}

\begin{interpretation}
Les deux couples donnent ici des conclusions \textbf{concordantes} :
\begin{itemize}[leftmargin=*]
  \item L'\textbf{analyseur A} est \textbf{précis et fidèle} : ses deux indicateurs de dispersion (écart type et IQR) sont faibles, et la moyenne est très proche de la cible 1,00 g/L.
  \item L'\textbf{analyseur B} est \textbf{moins fidèle} : ses indicateurs de dispersion sont 4 à 5 fois plus élevés. La moyenne (1,03 g/L) et la médiane (1,01 g/L) sont également plus éloignées de la cible.
\end{itemize}

\vspace{0.2cm}
\textbf{Conclusion :} l'analyseur A doit être retenu pour le contrôle qualité. L'analyseur B doit être \textbf{recalibré} ou \textbf{révisé}.
\end{interpretation}

\begin{attention}
\textbf{Que se passerait-il si la série B contenait un unique résultat aberrant (ex. : 5,00 g/L) ?}

\vspace{0.2cm}
La moyenne B passerait de 1,03 à environ 1,43 g/L et l'écart type exploserait à environ 1,2 g/L. En revanche, la médiane et l'IQR resteraient quasiment inchangés (voir Figure 4). Cela montre la \textbf{robustesse} du couple (médiane ; IQR) face aux valeurs extrêmes. C'est pourquoi les laboratoires utilisent souvent les deux couples \textbf{en parallèle} : la moyenne et l'écart type pour la fidélité, la médiane et l'IQR pour la description et la détection d'anomalies.
\end{attention}

% ============================================================
\newpage
% ============================================================

\section*{6. Synthèse à retenir}

\begin{consigne}
\textbf{À mémoriser :}

\vspace{0.2cm}
\begin{itemize}[leftmargin=*]
  \item \textbf{Écart type $\sigma$ ou $s$} : dispersion \textbf{moyenne} autour de la moyenne. Sensible aux valeurs extrêmes. Unité = celle de la variable. Utile pour la \textbf{fidélité}, les \textbf{cartes de contrôle}, les \textbf{tests statistiques}.

  \item \textbf{Écart interquartile IQR} $= Q_3 - Q_1$ : dispersion des \textbf{50\,\% centraux}. Robuste aux valeurs extrêmes. Unité = celle de la variable. Utile pour les \textbf{distributions asymétriques}, les \textbf{boîtes à moustaches}, la \textbf{comparaison de populations hétérogènes}.

  \item \textbf{Choix du couple :}
        \begin{itemize}
          \item Distribution \textbf{symétrique}, sans valeurs aberrantes → \textbf{(moyenne ; écart type)}.
          \item Distribution \textbf{asymétrique} ou avec valeurs aberrantes → \textbf{(médiane ; IQR)}.
          \item En cas de doute → \textbf{calculer les deux} et comparer $\bar{x}$ et $Me$.
        \end{itemize}

  \item \textbf{Complémentarité :} en pratique, on gagne souvent à présenter \textbf{les deux couples} (ou une boîte à moustaches enrichie de la moyenne) pour donner une vision complète de la série.
\end{itemize}
\end{consigne}

\vspace{0.4cm}

\section*{7. Exercices d'application}

\begin{exemple}
\textbf{Exercice 1.} On mesure la concentration d'un analyte (en mg/L) sur 9 échantillons :
\[
12{,}1 \quad 12{,}3 \quad 12{,}0 \quad 12{,}4 \quad 12{,}2 \quad 12{,}1 \quad 12{,}5 \quad 12{,}2 \quad 12{,}3
\]
\begin{enumerate}[leftmargin=*]
  \item Calculer la moyenne et l'écart type (corrigé).
  \item Calculer la médiane et l'écart interquartile.
  \item Quelle conclusion tirez-vous sur la fidélité de la méthode ?
\end{enumerate}
\lignes{5}
\end{exemple}

\begin{exemple}
\textbf{Exercice 2.} On mesure le temps de coagulation (en secondes) sur 10 patients :
\[
28 \quad 31 \quad 29 \quad 30 \quad 32 \quad 45 \quad 30 \quad 29 \quad 31 \quad 30
\]
\begin{enumerate}[leftmargin=*]
  \item Calculer la moyenne et l'écart type.
  \item Calculer la médiane et l'écart interquartile.
  \item Les deux couples donnent-ils la même vision de la série ? Lequel choisiriez-vous pour décrire cette série ? Justifier.
  \item Construire la boîte à moustaches correspondante.
\end{enumerate}
\lignes{6}
\end{exemple}

\begin{exemple}
\textbf{Exercice 3 (synthèse).} Un laboratoire souhaite valider une nouvelle méthode de dosage. Rédigez un court paragraphe (8 lignes) expliquant :
\begin{itemize}[leftmargin=*]
  \item quel couple d'indicateurs vous utiliseriez pour évaluer la \textbf{fidélité} de la méthode ;
  \item quel couple vous utiliseriez pour \textbf{décrire} la population de patients testés ;
  \item pourquoi il est utile de présenter \textbf{les deux couples} dans le rapport final.
\end{itemize}
\lignes{8}
\end{exemple}

\vfill

\begin{center}
\textit{Fiche méthode — BTS BioALC — 2\textsuperscript{ème} année — Couples (moyenne ; $\sigma$) et (médiane ; IQR)}\\
\textit{Document à conserver pour les révisions et les travaux pratiques.}
\end{center}

\end{document}